\section{横向全家桶与纵向垂直整合}

上一章讲分化，本章讲收敛。C 端有明显头部收敛趋势，ChatGPT 可能收敛掉很多产品。横向全家桶的例子是 ChatGPT，把 Chat、搜索、Coding、Agent、Workspace 都纳入同一入口；纵向垂直整合的例子是 Gemini，从 TPU 芯片、Gemini 模型、Agent 应用，到 Google Docs、Chrome、Android、YouTube，可以形成超级集成。

\begin{figure}[H]
\centering
\includegraphics[width=0.94\textwidth]{horizontal-vs-vertical.png}
\caption{横向全家桶 vs 纵向垂直整合：ChatGPT 横向套件，Gemini 纵向超级集成。自制概念图，依据 00:21:37--00:33:35 对谈内容整理。}
\end{figure}

\begin{knowledgebox}{读图：横向和纵向是两种吞噬方式}
横向全家桶把多个用户任务聚到一个入口；纵向垂直整合把芯片、模型、应用、浏览器、操作系统和内容生态串起来。前者吃用户心智，后者吃基础设施和默认入口。
\end{knowledgebox}

\subsection{ChatGPT 的横向产品壁垒}

ChatGPT 不只是一个模型壳，它拥有用户习惯、品牌心智、默认入口、历史对话和持续功能扩展。对于独立 AI 产品，最难的是：你做出的功能如果变成 ChatGPT 全家桶里的一个标签页，用户为什么还要单独使用你？

\begin{importantbox}{横向全家桶的杀伤力}
横向全家桶最可怕的地方不是单点功能都最好，而是它成为用户默认入口。默认入口可以把“够好”的功能迅速分发给海量用户，从而压缩创业产品的窗口期。
\end{importantbox}

\subsection{Gemini 的纵向整合}

Google 的优势在纵向整合：TPU、云、Gemini、搜索、广告、Docs、Chrome、Android、YouTube 都能协同。广密对 Google 的看法发生转变，原因也在这里：当 AI 产品最终需要搜索、广告、办公、浏览器和操作系统协同，Google 的老资产可能重新变成新优势。

\begin{knowledgebox}{Google 纵向整合的资产}
\begin{tabularx}{\textwidth}{p{0.22\textwidth}p{0.34\textwidth}X}
\toprule
资产 & 作用 & AI 时代价值 \\
\midrule
TPU/Cloud & 算力和训练基础设施 & 降低模型成本并保证供给。 \\
Search/Ads & 入口和商业化系统 & AI 搜索和广告平台可能重新融合。 \\
Chrome/Android & 浏览器和移动操作系统 & Agent 可以进入默认数字环境。 \\
Docs/Workspace & 办公场景和用户文件 & 企业和个人生产力入口。 \\
YouTube & 视频内容和多模态数据/场景 & 多模态理解和生成的应用场。 \\
\bottomrule
\end{tabularx}
\end{knowledgebox}

\subsection{本章小结}

横向全家桶和纵向垂直整合解释了为什么头部公司越来越难绕开。创业者要么成为全家桶里的关键能力，要么找到巨头不容易整合的场景。

